% Generated by scripts/assemble.py; corrections belong in source/corrections.json.
\examhead{98}{南京航空航天大学 2024 年数学分析试卷}{97}
% p098-b14; source PDF page 98
\blockanchor{p098-b14}
\question{1}解答如下问题:\par
% p098-b15; source PDF page 98
\blockanchor{p098-b15}
(1) 求极限 $\lim_{x\to 0}\frac{1}{x^4}\left[6 + x^2 -\frac{5}{6} x^4 -6\sqrt[3]{2 - \cos x}\right].$\par
% p098-b16; source PDF page 98
\blockanchor{p098-b16}
(2) 求定积分  $\int_{0}^{\frac{\pi}{4}}\frac{x}{1+\cos2x}dx.$\par
% p098-b17; source PDF page 98
\blockanchor{p098-b17}
\question{2}设 $f(x) = \frac{x + 2}{x + 1} \sin \frac{1}{x}, a > 0$ 为常数, 证明: $f(x)$ 在 $[a, +\infty)$ 上一致连续, 在 $(0, a)$ 上不一致连续.\par
% p098-b18; source PDF page 98
\par\Needspace{6\baselineskip}
\blockanchor{p098-b18}
\question{3}设 $f(x)$ 在 $[a, b]$ 上连续, 在 $(a, b)$ 内二阶可导, 证明: 对任意的 $x \in [a, b]$, 存在 $\xi \in [a, b]$, 使得\par
\nopagebreak[4]
% p098-b19; source PDF page 98
\blockanchor{p098-b19}
\displayformula{\frac {f (x) - f (a)}{x - a} - \frac {f (b) - f (a)}{b - a} = \frac {1}{2} f ^ {\prime \prime} (\xi) (x - b).}
% p098-b20; source PDF page 98
\blockanchor{p098-b20}
\question{4}判断反常积分 $\int_0^{+\infty} \frac{(\ln^2 x) \sin x}{x + 1} \, \mathrm{d}x$ 是绝对收敛, 条件收敛还是发散, 并证明.\par
% p098-b21; source PDF page 98
\blockanchor{p098-b21}
\question{5}解答如下问题:\par
% p098-b22; source PDF page 98
\par\Needspace{6\baselineskip}
\blockanchor{p098-b22}
(1) 设 $f(x), g(x)$ 在 $[a, b]$ 上可积, 证明 Cauchy-Schwarz 不等式:\par
\nopagebreak[4]
% p098-b23; source PDF page 98
\blockanchor{p098-b23}
\displayformula{\left(\int_ {a} ^ {b} f (x) g (x) \mathrm{d} x\right) ^ {2} \leq \int_ {a} ^ {b} f ^ {2} (x) \mathrm{d} x \int_ {a} ^ {b} g ^ {2} (x) \mathrm{d} x.}
% p098-b24; source PDF page 98
\par\Needspace{6\baselineskip}
\blockanchor{p098-b24}
(2) 设 $f(x)$ 在 $[a, b]$ 上有连续的导函数, 且 $f(a) = 0$, 证明:\par
\nopagebreak[4]
% p098-b25; source PDF page 98
\blockanchor{p098-b25}
\displayformula{\int_ {a} ^ {b} | f (x) f ^ {\prime} (x) | \mathrm{d} x \leq \frac {b - a}{2} \int_ {a} ^ {b} [ f ^ {\prime} (x) ] ^ {2} \mathrm{d} x.}
% p099-b00; source PDF page 99
\blockanchor{p099-b00}
\question{6}设 $a_1 = 2, a_{n+1} = \frac{1}{2}\left(a_n + \frac{1}{a_n}\right), n = 1, 2, \cdots$.\par
% p099-b01; source PDF page 99
\blockanchor{p099-b01}
(1) 证明 $\{a_{n}\}$ 收敛, 并求其极限.\par
% p099-b02; source PDF page 99
\blockanchor{p099-b02}
(2) 判断级数 $\sum_{n=1}^{\infty}\left(\frac{a_n}{a_{n+1}} - 1\right)$ 的敛散性并证明.\par
% p099-b03; source PDF page 99
\blockanchor{p099-b03}
\question{7}将 $f(x) = e^{x} \left(-\pi \leq x \leq \pi\right)$ 展开成傅里叶级数, 讨论级数的收敛性, 并求级数 $\sum_{n=0}^{\infty} \frac{1}{1 + n^2}$ 的和.\par
% p099-b04; source PDF page 99
\par\Needspace{6\baselineskip}
\blockanchor{p099-b04}
\question{8}给定方程组\par
\nopagebreak[4]
% p099-b05; source PDF page 99
\blockanchor{p099-b05}
\displayformula{\left\{ \begin{array}{l} x u ^ {2} + v = y ^ {3}; \\ 2 y u - x v ^ {3} = 4 x. \end{array} \right.}
% p099-b06; source PDF page 99
\blockanchor{p099-b06}
判断其在 $P_0(x_0, y_0, u_0, v_0) = (0, 1, 0, 1)$ 附近能否确定 $u, v$ 为 $x, y$ 的函数, 若能, 求出偏导数 $\frac{\partial u}{\partial y}, \frac{\partial v}{\partial y}, \frac{\partial^2 v}{\partial y^2}$ 在点 $(x_0, y_0) = (0, 1)$ 的值.\par
% p099-b07; source PDF page 99
\blockanchor{p099-b07}
\question{9}解答如下问题:\par
% p099-b08; source PDF page 99
\par\Needspace{6\baselineskip}
\blockanchor{p099-b08}
(1) 求曲面积分\par
\nopagebreak[4]
% p099-b09; source PDF page 99
\blockanchor{p099-b09}
\displayformula{\iint_ {S} \left(x ^ {2} + y ^ {2}\right) \mathrm{d} S.}
% p099-b10; source PDF page 99
\blockanchor{p099-b10}
其中 $S$ 为立体 $\sqrt{x^2 + y^2} \leq z \leq 1$ 的边界曲面.\par
% p099-b11; source PDF page 99
\par\Needspace{6\baselineskip}
\blockanchor{p099-b11}
(2) 求曲线积分\par
\nopagebreak[4]
% p099-b12; source PDF page 99
\blockanchor{p099-b12}
\displayformula{\oint_ {L} (y + \sin x) \mathrm{d} x + (z - e ^ {y}) \mathrm{d} y + (x + 1) \mathrm{d} z.}
% p099-b13; source PDF page 99
\blockanchor{p099-b13}
其中 $L$ 为球面 $x^{2} + y^{2} + z^{2} = R^{2}$ 与平面 $x + y + z = 0$ 的交线，从 $x$ 轴正向看去，$L$ 沿逆时针方向.\par
% p099-b14; source PDF page 99
\par\Needspace{6\baselineskip}
\blockanchor{p099-b14}
\question{10}设 $f(x)$ 在 $[0, +\infty)$ 上可导, 且 $f(0) = 0$. 求极限\par
\nopagebreak[4]
% p099-b15; source PDF page 99
\blockanchor{p099-b15}
\displayformula{\lim _ {t \rightarrow 0 ^ {+}} \frac {1}{\pi t ^ {4}} \iiint_ {V} f (\sqrt {x ^ {2} + y ^ {2} + z ^ {2}}) d x d y d z.}
% p099-b16; source PDF page 99
\blockanchor{p099-b16}
其中 $V: x^{2} + y^{2} + z^{2} \leq t^{2}$.\par
% p099-b17; source PDF page 99
\blockanchor{p099-b17}
\question{11}设 $\{f_n(x)\}$ 是在 $[a, b]$ 上连续的函数列, 且 $\{f_n(x)\}$ 在 $[a, b]$ 上一致收敛于 $f(x), f(x)$ 在 $[a, b]$ 上恒为负值. 证明: 当 $n$ 充分大时, $f_n(x)$ 在 $[a, b]$ 上也恒为负值, 且函数列 $\left\{\frac{1}{f_n(x)}\right\}$ 在 $[a, b]$ 上一致收敛于 $\frac{1}{f(x)}$.\par
% p099-b18; source PDF page 99
\blockanchor{p099-b18}
\question{12}设 $f(x)$ 在 $[a, b]$ 上可积, 且 $A \leq f(x) \leq B$, $g(u)$ 在 $[A, B]$ 上连续, 证明: $g(f(x))$ 在 $[a, b]$ 上可积.\par
